\MathBookSourcePage{97}
\setcounter{statement}{1}
\begin{Remark}\label{rem:stokes-dirichlet-variational-formulations}
问题 \eqref{eq:stokes-dirichlet-problem} 分别具有 $(Q)$ 型和 $(P)$ 型变分形式:

求一对
$(\mathbf{u},p)\in H^1(\Omega)^N\times L_0^2(\Omega)$, 使
\begin{equation}\label{eq:stokes-dirichlet-q-formulation}
\left\{
\begin{aligned}
a(\mathbf{u},\mathbf{v})
-(p,\operatorname{div}\mathbf{v})
&=\langle\mathbf{f},\mathbf{v}\rangle
&&\forall\mathbf{v}\in H_0^1(\Omega)^N,\\
(q,\operatorname{div}\mathbf{u})&=0
&&\forall q\in L_0^2(\Omega),\\
\mathbf{u}&=\mathbf{g}
&&\text{on }\Gamma,
\end{aligned}
\right.
\end{equation}
以及

求 $\mathbf{u}\in H^1(\Omega)^N$, 使
\begin{equation}\label{eq:stokes-dirichlet-p-formulation}
\left\{
\begin{aligned}
a(\mathbf{u},\mathbf{v})
&=\langle\mathbf{f},\mathbf{v}\rangle
&&\forall\mathbf{v}\in V,\\
\operatorname{div}\mathbf{u}&=0
&&\text{in }\Omega,\\
\mathbf{u}&=\mathbf{g}
&&\text{on }\Gamma,
\end{aligned}
\right.
\end{equation}
其中
\[
a(\mathbf{u},\mathbf{v})
=
\left\{
\begin{aligned}
\nu(\operatorname{grad}\mathbf{u},
\operatorname{grad}\mathbf{v}),\\
2\nu\sum_{i,j=1}^{N}
(D_{ij}(\mathbf{u}),D_{ij}(\mathbf{v})).
\end{aligned}
\right.
\]
\end{Remark}

\setcounter{statement}{2}
\begin{Remark}\label{rem:stokes-dirichlet-apriori-bound}
Corollary \ref{cor:abstract-v-elliptic-well-posedness} 给出估计
\[
\|\mathbf{u}\|_{1,\Omega}+\|p\|_{0,\Omega}
\le C(\|\mathbf{f}\|_{-1,\Omega}
+\|\mathbf{u}_0\|_{1,\Omega}),
\]
它对所有满足
$\operatorname{div}\mathbf{u}_0=0$,
$\mathbf{u}_0|_\Gamma=\mathbf{g}$ 的
$\mathbf{u}_0\in H^1(\Omega)^N$ 成立. 对 $\mathbf{u}_0$ 取下确界,
便得到
\[
\|\mathbf{u}\|_{1,\Omega}+\|p\|_{0,\Omega}
\le C(\|\mathbf{f}\|_{-1,\Omega}
+\|\mathbf{g}\|_{1/2,\Gamma}).
\]
\end{Remark}

问题 \eqref{eq:stokes-dirichlet-problem} 也可表述成鞍点问题.
采用上述记号, 置
\begin{equation}\label{eq:stokes-dirichlet-energy}
J(\mathbf{v})
=\frac12a(\mathbf{v},\mathbf{v})
-\langle\mathbf{f},\mathbf{v}\rangle,
\end{equation}
\begin{equation}\label{eq:stokes-dirichlet-lagrangian}
\mathcal{L}(\mathbf{v},q)
=J(\mathbf{v})-(q,\operatorname{div}\mathbf{v}).
\end{equation}

\setcounter{statement}{1}
\begin{Theorem}\label{thm:stokes-dirichlet-saddle-characterization}
在 Theorem \ref{thm:stokes-dirichlet-well-posedness} 的假设下,
\eqref{eq:stokes-dirichlet-problem} 的解 $(\mathbf{u},p)$ 由下式刻画:
\begin{equation}\label{eq:stokes-dirichlet-minmax}
\begin{aligned}
\mathcal{L}(\mathbf{u},p)
&=
\operatorname*{Min}_{\substack{
\mathbf{v}\in H^1(\Omega)^N\\
\mathbf{v}|_\Gamma=\mathbf{g}}}
\ \sup_{q\in L_0^2(\Omega)}
\mathcal{L}(\mathbf{v},q)\\
&=
\operatorname*{Max}_{q\in L_0^2(\Omega)}
\ \inf_{\substack{
\mathbf{v}\in H^1(\Omega)^N\\
\mathbf{v}|_\Gamma=\mathbf{g}}}
\mathcal{L}(\mathbf{v},q).
\end{aligned}
\end{equation}
\end{Theorem}
\begin{Proof}
采用 \eqref{eq:stokes-dirichlet-abstract-data} 中的记号.
由于形式 $a(\cdot,\cdot)$ 对称且在 $H_0^1(\Omega)^N$ 上椭圆,
形式 $b(\cdot,\cdot)$ 满足 inf-sup 条件
\eqref{eq:stokes-pressure-infsup}, 可以应用 Theorem
\ref{thm:abstract-saddle-equivalence}. 得到
$(\mathbf{u}-\mathbf{u}_0,p)$ 是 Lagrange 泛函
\[
\mathcal{L}_{\mathbf{u}_0}(\mathbf{v},q)
=\frac12a(\mathbf{v},\mathbf{v})
-\langle\mathbf{f},\mathbf{v}\rangle
+a(\mathbf{u}_0,\mathbf{v})
-(q,\operatorname{div}\mathbf{v})
\]
在乘积空间 $H_0^1(\Omega)^N\times L_0^2(\Omega)$ 上的唯一鞍点.

\MathBookSourcePage{98}
考虑到
$\operatorname{div}\mathbf{u}_0
=\operatorname{div}\mathbf{u}=0$, 这恰好意味着
\[
\mathcal{L}(\mathbf{u},q)
=\mathcal{L}(\mathbf{u},p)
\le
\mathcal{L}(\mathbf{v}+\mathbf{u}_0,p)
\]
对所有 $\mathbf{v}\in H_0^1(\Omega)^N$ 和所有
$q\in L_0^2(\Omega)$ 成立, 因此
$(\mathbf{u},p)$ 是泛函 $\mathcal{L}$ 在
$[\mathbf{u}_0+H_0^1(\Omega)^N]\times L_0^2(\Omega)$ 上的唯一鞍点.

现在利用
\[
\mathbf{u}_0+H_0^1(\Omega)^N
=\{\mathbf{v}\in H^1(\Omega)^N;
\ \mathbf{v}|_\Gamma=\mathbf{g}\},
\]
并像 Corollary \ref{cor:abstract-saddle-minmax} 的证明那样论证,
即可得到刻画 \eqref{eq:stokes-dirichlet-minmax}.
\end{Proof}

调整 Section \ref{sec:abstract-variational-problem} 中鞍点方法的论证,
可知 $\mathbf{u}$ 由下式刻画:
\begin{equation}\label{eq:stokes-dirichlet-constrained-energy}
J(\mathbf{u})
=
\inf_{\substack{
\mathbf{v}\in H^1(\Omega)^N\\
\operatorname{div}\mathbf{v}=0,\ 
\mathbf{v}|_\Gamma=\mathbf{g}}}
J(\mathbf{v}).
\end{equation}
另一方面, 对每个 $q\in L_0^2(\Omega)$, 关联边值问题
\begin{equation}\label{eq:stokes-dirichlet-dual-state}
\left\{
\begin{aligned}
-\nu\Delta\mathbf{u}(q)
&=\mathbf{f}-\operatorname{grad}q
&&\text{in }\Omega,\\
\mathbf{u}(q)&=\mathbf{g}
&&\text{on }\Gamma
\end{aligned}
\right.
\end{equation}
的解 $\mathbf{u}(q)\in H^1(\Omega)^N$. 然后置
\begin{equation}\label{eq:stokes-dirichlet-dual-functional}
K(q)=\frac12a(\mathbf{w}(q),\mathbf{w}(q)),
\end{equation}
其中
$\mathbf{w}(q)=\mathbf{u}(q)-\mathbf{u}_0$,
$\operatorname{div}\mathbf{u}_0=0$,
$\mathbf{u}_0|_\Gamma=\mathbf{g}$, 则
$p\in L_0^2(\Omega)$ 由
\begin{equation}\label{eq:stokes-dirichlet-dual-minimum}
K(p)=\operatorname*{Min}_{q\in L_0^2(\Omega)}K(q)
\end{equation}
刻画.

为了消去压力, 可以使用 Subsection
\ref{subsec:abstract-regularization-penalty} 中介绍的正则化方法或罚方法.
考虑问题
\begin{equation}\label{eq:stokes-dirichlet-penalty-system}
\left\{
\begin{aligned}
-\nu\Delta\mathbf{u}^\varepsilon
+\operatorname{grad}p^\varepsilon
&=\mathbf{f}
&&\text{in }\Omega,\\
p^\varepsilon
&=-\frac1\varepsilon
\operatorname{div}\mathbf{u}^\varepsilon
&&\text{in }\Omega,\\
\mathbf{u}^\varepsilon&=\mathbf{g}
&&\text{on }\Gamma.
\end{aligned}
\right.
\end{equation}
消去 $p^\varepsilon$, 得到关于 $\mathbf{u}^\varepsilon$
的等价二阶椭圆问题
\begin{equation}\label{eq:stokes-dirichlet-penalized-elliptic}
\left\{
\begin{aligned}
-\nu\Delta\mathbf{u}^\varepsilon
-\frac1\varepsilon
\operatorname{grad}\operatorname{div}\mathbf{u}^\varepsilon
&=\mathbf{f}
&&\text{in }\Omega,\\
\mathbf{u}^\varepsilon&=\mathbf{g}
&&\text{on }\Gamma.
\end{aligned}
\right.
\end{equation}

\MathBookSourcePage{99}
\setcounter{statement}{2}
\begin{Theorem}\label{thm:stokes-dirichlet-penalty-error}
假设 Theorem \ref{thm:stokes-dirichlet-well-posedness} 的条件成立.
则存在唯一一对函数
$(\mathbf{u}^\varepsilon,p^\varepsilon)
\in H^1(\Omega)^N\times L_0^2(\Omega)$,
它是方程 \eqref{eq:stokes-dirichlet-penalty-system} 的解. 此外有估计
\begin{equation}\label{eq:stokes-dirichlet-penalty-error-bound}
\|\mathbf{u}^\varepsilon-\mathbf{u}\|_{1,\Omega}
+\|p^\varepsilon-p\|_{0,\Omega}
\le
C\varepsilon
(\|\mathbf{f}\|_{-1,\Omega}
+\|\mathbf{g}\|_{1/2,\Gamma}),
\end{equation}
其中常数 $C$ 与 $\varepsilon,\mathbf{f},\mathbf{g}$ 无关.
\end{Theorem}
\begin{Proof}
与 Theorem \ref{thm:stokes-dirichlet-well-posedness} 一样,
容易得到问题 \eqref{eq:stokes-dirichlet-penalty-system}
的解对 $(\mathbf{u}^\varepsilon,p^\varepsilon)$ 的存在唯一性.
再考虑到差
$\mathbf{u}^\varepsilon-\mathbf{u}$ 在 $\Gamma$ 上为零,
Theorem \ref{thm:abstract-penalty-error} 的论证立即给出
\[
\|\mathbf{u}^\varepsilon-\mathbf{u}\|_{1,\Omega}
+\|p^\varepsilon-p\|_{0,\Omega}
\le C_1\varepsilon\|p\|_{0,\Omega},
\]
其中常数 $C_1$ 与 $\varepsilon$ 无关.
所以 \eqref{eq:stokes-dirichlet-penalty-error-bound}
由 Remark \ref{rem:stokes-dirichlet-apriori-bound} 得到.
\end{Proof}

类似地, 也可以应用 Theorem
\ref{thm:abstract-penalty-asymptotic-expansion}. 从 $p_0=p$ 开始,
问题
\begin{equation}\label{eq:stokes-dirichlet-asymptotic-recursion}
\left\{
\begin{aligned}
-\nu\Delta\mathbf{u}_n+\operatorname{grad}p_n
&=\mathbf{0}
&&\text{in }\Omega,\\
\operatorname{div}\mathbf{u}_n&=-p_{n-1}
&&\text{in }\Omega,\\
\mathbf{u}_n&=\mathbf{0}
&&\text{on }\Gamma
\end{aligned}
\right.
\end{equation}
通过归纳唯一确定序列
$(\mathbf{u}_n,p_n)\in H_0^1(\Omega)^N\times L_0^2(\Omega)$.
此外, 对所有 $M\ge1$ 以及足够小的 $\varepsilon$, 有
\begin{equation}\label{eq:stokes-dirichlet-asymptotic-error}
\begin{aligned}
&\left\|
\mathbf{u}^\varepsilon-\mathbf{u}
-\sum_{n=1}^{M}\varepsilon^n\mathbf{u}_n
\right\|_{1,\Omega}
+\left\|
p^\varepsilon-p
-\sum_{n=1}^{M}\varepsilon^np_n
\right\|_{0,\Omega}\\
&\qquad\le
C_M\varepsilon^{M+1}
(\|\mathbf{f}\|_{-1,\Omega}
+\|\mathbf{g}\|_{1/2,\Gamma}),
\end{aligned}
\end{equation}
其中常数 $C_M$ 与 $\varepsilon,\mathbf{f},\mathbf{g}$ 无关.

\setcounter{statement}{3}
\begin{Remark}\label{rem:stokes-dirichlet-strain-penalty}
若选择
\[
a(\mathbf{u},\mathbf{v})
=2\nu\sum_{i,j=1}^{N}
\int_\Omega
D_{ij}(\mathbf{u})D_{ij}(\mathbf{v})\,\mathrm{d}x,
\]
则得到另一个正则化问题:
\[
\left\{
\begin{aligned}
-2\nu\sum_{j=1}^{N}
\frac{\partial D_{ij}(\mathbf{u}^\varepsilon)}
{\partial x_j}
+\frac{\partial p^\varepsilon}{\partial x_i}
&=f_i
&&\text{in }\Omega,\quad 1\le i\le N,\\
p^\varepsilon
&=-\frac1\varepsilon
\operatorname{div}\mathbf{u}^\varepsilon
&&\text{in }\Omega,\\
\mathbf{u}^\varepsilon&=\mathbf{g}
&&\text{on }\Gamma.
\end{aligned}
\right.
\]
消去 $p^\varepsilon$, 得到罚问题
\[
\left\{
\begin{aligned}
-\nu\Delta\mathbf{u}^\varepsilon
-(\nu+1/\varepsilon)
\operatorname{grad}\operatorname{div}\mathbf{u}^\varepsilon
&=\mathbf{f}
&&\text{in }\Omega,\\
\mathbf{u}^\varepsilon&=\mathbf{g}
&&\text{on }\Gamma.
\end{aligned}
\right.
\]

\MathBookSourcePage{100}
为应用 Theorem \ref{thm:abstract-penalty-error}, 这里需要检验假设
\eqref{eq:abstract-combined-ellipticity}, 即
\begin{equation}\label{eq:stokes-strain-penalty-coercivity}
2\nu\sum_{i,j=1}^{N}
\|D_{ij}(\mathbf{v})\|_{0,\Omega}^2
+\|\operatorname{div}\mathbf{v}\|_{0,\Omega}^2
\ge
\alpha_0|\mathbf{v}|_{1,\Omega}^2
\quad
\forall\mathbf{v}\in H_0^1(\Omega)^N,
\qquad \alpha_0>0.
\end{equation}
但 Remark \ref{rem:stokes-pressure-quotient-space} 已经证明
\[
\begin{aligned}
\sum_{i,j=1}^{N}
(D_{ij}(\mathbf{u}),D_{ij}(\boldsymbol{\phi}))
&=-\frac12\{
\langle\Delta\mathbf{u},\boldsymbol{\phi}\rangle
+\langle
\operatorname{grad}(\operatorname{div}\mathbf{u}),
\boldsymbol{\phi}
\rangle\}\\
&=\frac12\{
(\operatorname{grad}\mathbf{u},
\operatorname{grad}\boldsymbol{\phi})
+(\operatorname{div}\mathbf{u},
\operatorname{div}\boldsymbol{\phi})
\}
\end{aligned}
\]
对所有 $\mathbf{u}\in H^1(\Omega)^N$ 和所有
$\boldsymbol{\phi}\in\mathcal{D}(\Omega)^N$ 成立. 因而
\begin{equation}\label{eq:stokes-strain-korn-h0-identity}
2\nu\sum_{i,j=1}^{N}
\|D_{ij}(\mathbf{v})\|_{0,\Omega}^2
=
\nu\{
|\mathbf{v}|_{1,\Omega}^2
+\|\operatorname{div}\mathbf{v}\|_{0,\Omega}^2
\}.
\end{equation}
所以在这种情况下, Theorem
\ref{thm:stokes-dirichlet-penalty-error} 的对应结论也成立.
\end{Remark}

\setcounter{statement}{4}
\begin{Remark}\label{rem:stokes-korn-inequality}
等式 \eqref{eq:stokes-strain-korn-h0-identity} 还与经典的
Korn 不等式有关:
\begin{equation}\tag{\ref*{eq:stokes-strain-korn-h0-identity}$'$}
\label{eq:stokes-korn-inequality}
\sum_{i,j=1}^{N}
\|D_{ij}(\mathbf{v})\|_{0,\Omega}^2
+\|\mathbf{v}\|_{0,\Omega}^2
\ge
\alpha_1\|\mathbf{v}\|_{1,\Omega}^2
\qquad
\forall\mathbf{v}\in H^1(\Omega)^N,
\qquad \alpha_1>0.
\end{equation}
为完整起见, 下面给出 Duvaut \& Lions [26]
针对 $\mathbb{R}^N$ 中有界 Lipschitz 连续开集给出的
\eqref{eq:stokes-korn-inequality} 的简要证明.

注意, 在形式上有恒等式
\[
\frac{\partial^2v_i}{\partial x_j\partial x_k}
=
\frac{\partial D_{ik}(\mathbf{v})}{\partial x_j}
+\frac{\partial D_{ij}(\mathbf{v})}{\partial x_k}
-\frac{\partial D_{jk}(\mathbf{v})}{\partial x_i},
\qquad
1\le i,j,k\le N.
\]
因此, 对 $\mathbf{v}\in H^1(\Omega)^N$, 有
\[
\begin{aligned}
\left\|
\operatorname{grad}\left(
\frac{\partial v_i}{\partial x_j}
\right)
\right\|_{-1,\Omega}
&\le
C_1\sum_{i,j=1}^{N}
\|\operatorname{grad}D_{ij}(\mathbf{v})\|_{-1,\Omega}\\
&\le
C_2\sum_{i,j=1}^{N}
\|D_{ij}(\mathbf{v})\|_{0,\Omega}.
\end{aligned}
\]
所以对 $\operatorname{grad}v_i$ 应用 Theorem
\ref{thm:necas-gradient-negative-norm},
立即得到 \eqref{eq:stokes-korn-inequality}.
\end{Remark}

现在转向求解问题 \eqref{eq:stokes-dirichlet-problem} 的梯度方法.
仍取
\[
c(p,q)=(p,q),
\]
于是
\[
a_r(\mathbf{u},\mathbf{v})
=\nu(\operatorname{grad}\mathbf{u},
\operatorname{grad}\mathbf{v})
+r(\operatorname{div}\mathbf{u},
\operatorname{div}\mathbf{v}),
\]
它在 $H_0^1(\Omega)^N$ 上显然是椭圆的. 因此, 用下式定义
$\mathbf{u}_r(q)$:

\MathBookSourcePage{101}
\[
\left\{
\begin{aligned}
a_r(\mathbf{u}_r(q),\mathbf{v})
&=\langle\mathbf{f},\mathbf{v}\rangle
+(\operatorname{div}\mathbf{v},q)
&&\forall\mathbf{v}\in H_0^1(\Omega)^N,\\
\mathbf{u}_r(q)&=\mathbf{g}
&&\text{on }\Gamma,
\end{aligned}
\right.
\]
由 \eqref{eq:stokes-dirichlet-dual-functional} 得
\[
\begin{aligned}
a_r(D\mathbf{u}_r\cdot\mu,\mathbf{v})
&=(\operatorname{div}\mathbf{v},\mu)
&&\forall\mathbf{v}\in H_0^1(\Omega)^N,\\
D\mathbf{u}_r\cdot\mu&=\mathbf{0}
&&\text{on }\Gamma,\\
DK_r(q)&=\operatorname{div}\mathbf{u}_r(q),\\
D^2K_r\cdot\mu&=
\operatorname{div}(D\mathbf{u}_r\cdot\mu).
\end{aligned}
\]
所以带最优参数的简单梯度算法如下.

$1^\circ)$ 给定 $p^0\in L_0^2(\Omega)$, 求解非齐次椭圆边值问题
\[
\left\{
\begin{aligned}
-(\nu\Delta+r\operatorname{grad}\operatorname{div})
\mathbf{u}^0
&=\mathbf{f}-\operatorname{grad}p^0
&&\text{in }\Omega,\\
\mathbf{u}^0&=\mathbf{g}
&&\text{on }\Gamma.
\end{aligned}
\right.
\]

$2^\circ)$ 对 $m\ge0$, 已知
$(\mathbf{u}^m,p^m)
\in H^1(\Omega)^N\times L_0^2(\Omega)$ 且
$\mathbf{u}^m|_\Gamma=\mathbf{g}$, 求解齐次问题
\[
\left\{
\begin{aligned}
-(\nu\Delta+r\operatorname{grad}\operatorname{div})
\mathbf{z}^m
&=\operatorname{grad}\operatorname{div}\mathbf{u}^m
&&\text{in }\Omega,\\
\mathbf{z}^m&=\mathbf{0}
&&\text{on }\Gamma;
\end{aligned}
\right.
\]
然后由
\[
\begin{aligned}
\rho^m
&=-\frac{\|\operatorname{div}\mathbf{u}^m\|_{0,\Omega}^2}
{(\operatorname{div}\mathbf{z}^m,
\operatorname{div}\mathbf{u}^m)},\\
p^{m+1}
&=p^m-\rho^m\operatorname{div}\mathbf{u}^m,\\
\mathbf{u}^{m+1}
&=\mathbf{u}^m+\rho^m\mathbf{z}^m
\end{aligned}
\]
计算 $\rho^m\in\mathbb{R}$ 及下一对
$(\mathbf{u}^{m+1},p^{m+1})
\in H^1(\Omega)^N\times L_0^2(\Omega)$.

容易检验 Theorem \ref{thm:abstract-general-descent-convergence}
的结论成立, 即上述格式总是收敛.
当 $\rho^m$ 不是最优参数时, 只要
\[
0<\inf_m\rho^m\le\sup_m\rho^m<2\widetilde{\alpha}_r,
\]
该结论仍然成立. 简单计算表明
\[
\widetilde{\alpha}_r\le r+\nu/N.
\]

共轭梯度算法具有相同的起始步骤 $1^\circ$, 而步骤
$2^\circ$ 替换为:

$2^\circ)$ 对 $m\ge0$, 已知
$(\mathbf{u}^m,p^m)\in H^1(\Omega)^N\times L_0^2(\Omega)$ 且
$\mathbf{u}^m|_\Gamma=\mathbf{g}$, 由
\[
\left.
\begin{aligned}
\sigma^m
&=\frac{\|\operatorname{div}\mathbf{u}^m\|_{0,\Omega}^2}
{\|\operatorname{div}\mathbf{u}^{m-1}\|_{0,\Omega}^2},\\
\omega^m
&=\operatorname{div}\mathbf{u}^m+\sigma^m\omega^{m-1}
\end{aligned}
\right\}
\quad\text{only if }m\ge1,
\qquad
\omega^0=\operatorname{div}\mathbf{u}^0\quad\text{otherwise},
\]
计算 $\sigma^m\in\mathbb{R}$ 和
$\omega^m\in L_0^2(\Omega)$.
